\chapter{Funciones booleanas}

\section{Estudio de las funciones booleanas sobre álgebras de Boole finitas}

Estudiaremos sólo las funciones $f: B^{n} \rightarrow B$ dónde $B = B_{2} = \{0, 1\}.$ Esto es $f: B_{2}^{n} \rightarrow B_{2}.$

\subsection{Formas normales}
Toda función $f: B_{2}^{n} \rightarrow B_{2}$ se puede poner en la forma:
\[
f(x_{1}, x_{2}, \dots, x_{n}) = \sum_{k=1}^{m < 2^{n}} \prod_{l=1}^{n} x_{l}
\]

\subsection{Número de funciones posibles}
En general, el número de funciones de un conjunto $C$ de cardinal $n_{C} \in \mathbb{N}$ en un conjunto $D$ de cardinal $n_{D} \in \mathbb{N}$ será $n_{D}^{n_{C}}.$ Así, una función de $n$ variables en $C$ con valores en $D$ será $n_{D}^{(n_{C}^{n})}.$ Para el caso en que $C = D = B_{2}$ que toma como argumento $n$ variables, obtenemos que el número de funciones será de $2^{(2^{n})}.$ 

En la siguiente tabla vemos cómo crece esta cantidad:

\begin{table}[h]
\centering
\begin{tabular}{lll}
\hline
\textbf{N variables} & \textbf{Longitud tabla ($2^N$)} & \textbf{Número de funciones distintas: $2^{(2^{N})}$} \\
\hline
0 & 1 & 2 \\
1 & 2 & 4 \\
2 & 4 & 16 \\
3 & 8 & 256 \\
4 & 16 & 65536 \\
5 & 32 & 4294967296 \\
6 & 64 & 18446744073709551616 \\
7 & 128 & 340282366920938463463374607431768211456 \\
\hline
\end{tabular}
\end{table}

\section{Demostración por Inducción del número de funciones}

El punto anterior es fácil de probar:

\begin{enumerate}
    \item Para el caso de $N = 0$ y de $N = 1$ es fácil probar (por enumeración) la validez de la fórmula. Más tarde mostraremos tablas de todas las funciones hasta $N = 2$ inclusive.
    \item Para el caso general, la Hipótesis de Inducción (HI) será:
    \[
    [\mathrm{HI}] \quad \forall N \in (\mathbb{N} \cup \{0\}), \ 0 \leq N \leq n - 1 \implies 2^{(2^{N})} \text{ es el cardinal buscado.}
    \]
    \item Veremos si para el caso $N = n$ se sigue cumpliendo la fórmula anterior. Pero esto es claro: al añadir una variable en el argumento tendremos todas las funciones del caso $N = n - 1$ ($2^{n - 1}$) para el valor $0$ de la nueva variable y otros $2^{n - 1}$ para el valor $1$ de la nueva variable, y no quedan otros casos. Las funciones totales para $N = n$ serán:
    \[
    \mathrm{card} \left( \left\{ f: B_{2}^{n - 1} \rightarrow B_{2} \right\} \times \left\{ f: B_{2}^{n - 1} \rightarrow B_{2} \right\} \right) =
    \]
    \[
    = \mathrm{card} \left( \left\{ f: B_{2}^{n - 1} \rightarrow B_{2} \right\} \right) \cdot \mathrm{card} \left( \left\{ f: B_{2}^{n - 1} \rightarrow B_{2} \right\} \right) =
    \]
    \[
    = 2^{(2^{n - 1})} \cdot 2^{(2^{n - 1})} = 2^{(2^{n - 1} + 2^{n - 1})} = 2^{2 \cdot (2^{n - 1})} = 2^{(2^{n})}
    \]
\end{enumerate}

Y así queda establecida la fórmula.

\section{Minitérminos y Maxitérminos}

Como se ve en el punto anterior, el crecimiento es desmesurado al compararlo al crecimiento lineal de los argumentos. En un futuro, cuando intentemos hacer reducciones de expresiones booleanas, este crecimiento nos impedirá construir métodos eficaces para resolver las minimizaciones.

Aunque hemos visto que podemos poner las expresiones booleanas en los conjuntos de operadores $\{ \{+, \overline{\quad}\}, \{\cdot, \overline{\quad}\}, \{\uparrow\}, \{\downarrow\}, \{\oplus, \cdot\}, \{\odot, +\} \},$ por comprensibilidad y para una lectura normal se utilizan frecuentemente los dos primeros conjuntos unidos, pudiendo variarse bien el orden de los operadores binarios: $\{ \{+, \cdot, \overline{\quad}\}, \{\cdot, +, \overline{\quad}\} \}.$ 

El primer conjunto $\{+, \cdot, \overline{\quad}\}$ será el que estudiemos por defecto, el segundo se tratará con una simetría de dualidad (hay que tener algunos cuidados). El conjunto elegido de operaciones se llamará desarrollo en sumas de productos de términos simples (una variable, o su negada, o una constante). También se llamará desarrollo por minitérminos o SOP (inglés) o SdP. El segundo será el desarrollo en producto de sumas de términos simples. También se llamará desarrollo por maxitérminos o POS (inglés) o PdS.

\subsection{Desarrollo por minitérminos}
En el desarrollo por minitérminos expresamos sólo los términos de la expresión en que la función tiene como valor $1.$ Veamos:

\[
f_{B_{2}}(x_{1}, x_{2}, \dots, x_{n}) = \sum_{(i_{1}, i_{2}, \dots, i_{n}) = (0, \dots, 0)}^{(1, \dots, 1)} \left( \sigma_{1}^{(i_{1}, i_{2}, \dots, i_{n})}(i_{1}) \cdot \sigma_{2}^{(i_{1}, i_{2}, \dots, i_{n})}(i_{2}) \cdot \dots \cdot \sigma_{n}^{(i_{1}, i_{2}, \dots, i_{n})}(i_{n}) \right)
\]

dónde $\sigma_{k}^{(i_{1}, i_{2}, \dots, i_{n})}(i_{k}) \in \{i_{k}, \overline{i_{k}}\}$ y $0 \leq k \leq n.$

Para más sencillez:

\[
f_{B_{2}}(x) = f_{B_{2}}(x_{1}, x_{2}, \dots, x_{n}) = \sum_{\iota \in B_{2}^{n}} \left( \sigma_{1}^{\iota}(\iota_{1}) \cdot \sigma_{2}^{\iota}(\iota_{2}) \cdot \dots \cdot \sigma_{n}^{\iota}(\iota_{n}) \cdot f_{B_{2}}(\iota) \right) = \sum_{\iota \in B_{2}^{n}} (\sigma^{\iota} \cdot f_{B_{2}}(\iota))
\]
dónde $\sigma_{k}^{\iota}(\iota_{k}) \in \{\iota_{k}, \overline{\iota_{k}}\},$ $0 \leq k \leq n$ y $\sigma^{\iota}$ es un minitérmino.

Cuándo la sigma (el minitérmino) es en todos los casos (para todas las iotas) completo (es un producto de $n$-términos simples), el valor de la función en esa iota concreta indica si el minitérmino aparece o no.

\subsection{Desarrollo por maxitérminos}
En el desarrollo por maxitérminos de una función de $n$ variables, los maxitérminos son sumatorios de $n$-términos simples. Así termina consistiendo la función en un producto de maxitérminos. Los maxitérminos indican un $0$ de la función.

\begin{quote}\textbf{Teorema (Universalidad de las Formas Canónicas (Expansión de Shannon)):}\label{shannon_expansion}
Toda función booleana $f: \mathbb{B}^n \to \mathbb{B}$ puede expresarse de manera única (salvo conmutatividad) de dos formas canónicas duales:
\begin{enumerate}
    \item \textbf{Suma de Productos (SOP) / Minitérminos:} $f(x_1, \dots, x_n) = \sum_{\iota \in \mathbb{B}^n} f(\iota) \cdot \sigma^\iota$
    \item \textbf{Producto de Sumas (POS) / Maxitérminos:} $f(x_1, \dots, x_n) = \prod_{\iota \in \mathbb{B}^n} (f(\iota) + \Pi^\iota)$
\end{enumerate}
donde $\sigma^\iota$ es el minitérmino (producto en el que todas las variables aparecen afirmadas si $\iota_k=1$ o negadas si $\iota_k=0$), y $\Pi^\iota$ es el maxitérmino dual.
\end{quote}
\begin{proof}
Demostraremos la forma de minitérminos (SOP). Sea $\vec{x} \in \mathbb{B}^n$ un vector de entrada arbitrario. Por la construcción de los minitérminos, el término $\sigma^\iota(\vec{x}) = 1$ si y solo si el vector de entrada $\vec{x}$ es exactamente idéntico al vector constante $\iota.$ Para cualquier otro vector $\iota' \ne \vec{x},$ se cumple $\sigma^{\iota'}(\vec{x}) = 0.$

Al evaluar la expresión completa $\sum_{\iota \in \mathbb{B}^n} f(\iota) \cdot \sigma^\iota(\vec{x}),$ todos los sumandos en los que $\iota \ne \vec{x}$ se anulan, ya que $\sigma^\iota(\vec{x}) = 0 \implies f(\iota) \cdot 0 = 0.$ El único sumando que sobrevive a la suma (OR) es aquel en el que $\iota = \vec{x},$ para el cual $\sigma^{\vec{x}}(\vec{x}) = 1.$ Por lo tanto, la suma se colapsa a:
\[ f(\vec{x}) \cdot 1 = f(\vec{x}) \]
Esto demuestra que la expresión reproduce exactamente la tabla de verdad de la función $f.$ La demostración para los maxitérminos (POS) es idéntica por el Principio de Dualidad, donde $\Pi^\iota(\vec{x}) = 0$ si y solo si $\vec{x} = \iota,$ colapsando el producto (AND) al valor $f(\vec{x}) + 0 = f(\vec{x}).$
\end{proof}

\subsection{Transformación a Compuertas XOR y AND (Zhegalkin)}

Hemos visto en el capítulo dedicado al Sistema de Zhegalkin (ver Capítulo \ref{capitulo:zhegalkin}) que el anillo booleano $\langle \mathbb{B}, \oplus, \cdot \rangle$ forma un conjunto funcionalmente completo. Surge la pregunta práctica: dada una función expresada en su forma canónica de minitérminos (con 2 capas lógicas de AND $\to$ OR), ¿cómo podemos transformarla sistemáticamente a un circuito equivalente usando exclusivamente puertas XOR y AND?

\begin{quote}\textbf{Teorema (Conversión de Minitérminos a Polinomio de Zhegalkin):}\label{minterms_to_zhegalkin}
Toda expresión canónica en Suma de Productos (SOP) puede transformarse directamente en un Polinomio de Zhegalkin (XOR-AND) aplicando dos reglas:
\begin{enumerate}
    \item Reemplazar directamente todas las sumas lógicas (OR, $+$) por sumas exclusivas (XOR, $\oplus$).
    \item Reemplazar cada variable negada $\overline{x_k}$ por $(x_k \oplus 1)$ y aplicar la distributividad del producto sobre la suma exclusiva.
\end{enumerate}
\end{quote}
\begin{proof}
Para demostrar la regla 1, recordamos la relación fundamental entre OR y XOR:
\[ A + B = A \oplus B \oplus (A \cdot B) \]
En una expansión canónica por minitérminos, cualquier par de minitérminos distintos $\sigma^i$ y $\sigma^j$ ($i \ne j$) son \textbf{mutuamente excluyentes}. Esto significa que nunca pueden valer $1$ simultáneamente para la misma entrada, por lo que su producto lógico es siempre falso: $\sigma^i \cdot \sigma^j = 0.$
Sustituyendo esto en la relación anterior, obtenemos:
\[ \sigma^i + \sigma^j = \sigma^i \oplus \sigma^j \oplus 0 = \sigma^i \oplus \sigma^j \]
Por inducción, una suma lógica de cualquier cantidad de minitérminos mutuamente excluyentes es estrictamente equivalente a su suma exclusiva. Así, la capa OR exterior puede ser sustituida por una capa XOR sin alterar la función matemática.

Para demostrar la regla 2, sabemos por la definición de la negación en el Anillo de Zhegalkin que $\overline{x} = x \oplus 1.$ Al realizar esta sustitución en los literales de cada minitérmino, obtenemos una expresión que solo contiene multiplicaciones ($\cdot$), sumas exclusivas ($\oplus$) y la constante $1.$ Dado que el operador AND distribuye sobre el XOR ($A \cdot (B \oplus C) = (A \cdot B) \oplus (A \cdot C)$), podemos expandir algebraicamente todos los paréntesis para obtener un polinomio multilineal puro compuesto únicamente por sumas exclusivas de productos no negados. Este resultado es el Polinomio de Zhegalkin único de la función.
\end{proof}

Además de expresar las funciones por cadenas de símbolos que constituyen un término, existe una posibilidad de expresar estas funciones por tablas lineales o por cuadros (tablas bidimensionales). Para cada combinación de valores booleanos a la entrada de una función obtenemos un valor booleano de salida.

Para que las funciones booleanas representen algo de interés para la ingeniería, lo primero que debemos tener es una forma de representar la información que queremos procesar. ¿Cómo representamos un número?. ¿Cómo una letra?. Haremos un alto en la exposición de funciones booleanas, para detallar más esta pregunta, poder responderla y así ver para qué estamos viendo las álgebras de Boole.
